\section{Agent Systems and Tasks}
\label{sec:methods}

\subsection{System Definition}

Building on multi-agent system formalism \citep{zhu2025establishing, guo2024large}, an \textbf{agent system} $\mathcal{S} = (A, E, C, \Omega)$ consists of a set of agents $A = \{a_1, \ldots, a_n\}$ (where $n \geq 1$), a shared environment $E$, a communication topology $C$, and an orchestration policy $\Omega$. When $|A| = 1$, we refer to this as a Single-Agent System (SAS); when $|A| > 1$, a Multi-Agent System (MAS). Each agent $a_i$ perceives, reasons, and acts within the shared environment via iterative feedback.


Formally, each agent $a_i$ is defined as a tuple $S_i=(\Phi_i,\mathcal{A}_i,M_i,\pi_i)$, where:
\begin{itemize}
    \item $\Phi_i$ is the reasoning policy (typically an LLM)
    \item $\mathcal{A}_i = \{\text{ToolCall}(t, \theta) : t \in \mathcal{T}, \theta \in \Theta_t\}$ is the action space consisting of tool usage, where $\mathcal{T}$ is the set of available tools (e.g., web search, code execution) and $\Theta_t$ represents valid parameter configurations for tool $t$
    \item $M_i$ is the internal memory
    \item $\pi_i: \mathcal{H} \rightarrow \mathcal{A}_i$ is the decision function mapping observation histories to actions
\end{itemize}

The observation history space $\mathcal{H}$ contains sequences of action-observation pairs. The decision function $\pi_i$ is instantiated by the reasoning policy $\Phi_i$ (the LLM): given a history $h_{i,t}$, the LLM generates a reasoning trace and selects the next action. 

For instance, a history $h_{i,t} = [(\text{``search(query=`pandas')''}, \text{``Found 5 files''}), ...]$ is processed by $\Phi_i$ to produce the next tool call $\alpha_{i,t+1}$.

At timestep $t$, agent $a_i$ selects an action $\alpha_{i,t} \in \mathcal{A}_i$ according to:
\[
\alpha_{i,t}=\pi_i(h_{i,t}), \quad 
o_{i,t}=E(\alpha_{i,t}), \quad 
h_{i,t+1}=f_i(h_{i,t},\alpha_{i,t},o_{i,t}),
\]
where $E$ denotes the environment and $h_{i,0} = \{s_0\}$ contains the initial task specification. The history update function $f_i: \mathcal{H} \times \mathcal{A}_i \times \mathcal{O} \rightarrow \mathcal{H}$ appends the new action-observation pair to the agent's history: $h_{i,t+1} = f_i(h_{i,t}, \alpha_{i,t}, o_{i,t}) = h_{i,t} \oplus (\alpha_{i,t}, o_{i,t})$, subject to context window truncation when $|h_{i,t+1}| > \text{MAX\_TOKENS}$. This update mechanism applies uniformly to both SAS and MAS configurations. Communication between agents occurs through explicit message passing in the orchestration layer.

\vspace{0.4em}
\paragraph{Single-Agent System (SAS).}
A \emph{Single-Agent System} contains one reasoning locus ($|A|=1$ where $A$ is the agent set). All perception, reasoning, and action occur within a single sequential loop, producing computational complexity $O(k)$ where $k$ is the number of reasoning iterations. SAS has zero communication overhead and minimal memory $O(k)$, but limited capacity for decomposition or verification.

\vspace{0.4em}
\paragraph{Multi-Agent System (MAS).}

A \emph{Multi-Agent System} is an agent system $\mathcal{S}$ with $|A| > 1$, where agents interact through communication topology $C$ and orchestration policy $\Omega$.

Communication topology $C$ defines information flow patterns between agents:
\begin{itemize}
    \item \textbf{Independent}: $C = \{(a_i, a_{\text{agg}}) : \forall i\}$ (agent-to-aggregator only, no peer communication)
    \item \textbf{Centralized}: $C = \{(a_{\text{orch}}, a_i) : \forall i\}$ (orchestrator-to-agents only)  
    \item \textbf{Decentralized}: $C = \{(a_i, a_j) : \forall i,j, i \neq j\}$ (all-to-all topology)
    \item \textbf{Hybrid}: $C = C_{\text{centralized}} \cup C_{\text{peer}}$ (orchestrator plus limited peer-to-peer)
\end{itemize}

The orchestrator $\Omega$ (when present) determines: (i) how sub-agent outputs are aggregated (e.g., majority voting, weighted synthesis), (ii) whether the orchestrator can override sub-agent decisions, (iii) whether memory persists across coordination rounds, and (iv) termination conditions based on consensus or quality thresholds.

MAS architectures vary by how information and control propagate among agents, creating distinct trade-offs between computation, coordination, and parallelization. Table \ref{tab:agent_methods_complexity} formalizes these trade-offs using asymptotic notations over \textit{LLM calls}, \textit{sequential depth}, \textit{communication overhead}, and \textit{memory complexity}. We selected these five architectures to form a \textbf{structural ablation of coordination mechanisms}:

\begin{itemize}
    \item \textbf{Independent} isolates the effect of parallelism (ensemble) without communication.
    \item \textbf{Decentralized} introduces peer-to-peer information fusion without hierarchy.
    \item \textbf{Centralized} introduces hierarchical verification and bottleneck control.
    \item \textbf{Hybrid} examines the combination of hierarchy and lateral flexibility.
\end{itemize}

This design allows us to systematically attribute performance gains to specific coordination mechanics rather than generic ``multi-agent'' effects. Specific configurations include:
\begin{itemize}
    \item \textbf{Independent MAS:} $A = \{a_1, \ldots, a_n\}$, $\mathcal{C} = \{(a_i, a_{\text{agg}})\}$, $\Omega = \texttt{synthesis\_only}$. The \texttt{synthesis\_only} policy concatenates sub-agent outputs without cross-validation or majority voting; the aggregator performs no analytical comparison of responses, ensuring that any performance differences arise purely from parallel exploration rather than error correction. This achieves maximal parallelization but minimal coordination, suitable for ensemble-style reasoning.
    \item \textbf{Centralized MAS}: $A = \{a_{\text{orch}}, a_1,\ldots,a_n\}$, $C = \{(a_{\text{orch}}, a_i) : \forall i\}$, $\Omega = \text{hierarchical}$. A single orchestrator coordinates $r$ rounds across $n$ sub-agents ($O(rnk)$). Sequential depth equals $r$ while parallelization factor remains $n$. This design stabilizes reasoning but creates a bottleneck at the orchestrator.
    \item \textbf{Decentralized MAS}: $A = \{a_1,\ldots,a_n\}$, $C = \{(a_i, a_j) : \forall i,j, i \neq j\}$, $\Omega = \text{consensus}$. Agents communicate in $d$ sequential debate rounds ($O(dnk)$). Memory complexity is $O(dnk)$ as each agent stores its own debate history. This enables consensus formation through peer-to-peer discussion.
    \item \textbf{Hybrid MAS}: $A = \{a_{\text{orch}}, a_1,\ldots,a_n\}$, $C = \text{star} + \text{peer edges}$, $\Omega = \text{hierarchical} + \text{lateral}$. Combines orchestrated hierarchy with limited peer communication ($O((r+p) \cdot n \cdot k)$ where $p$ is the number of peer rounds). This inherits orchestrator control while enabling lateral exchange between agents.
\end{itemize}

\vspace{0.4em}
\paragraph{Communication vs. Coordination.}
We distinguish \emph{communication} (message passing between agents) from \emph{coordination} (strategic direction of agent activities). In centralized systems, coordination occurs through the orchestrator's task decomposition and progress monitoring, while communication involves passing findings between orchestrator and workers. In decentralized systems, communication and coordination are intertwined through debate rounds where agents both exchange information and collectively steer problem-solving direction.

Thus, SAS represents the minimal unit of agentic computation ($O(k)$), while MAS configurations explore the scaling frontier of coordination complexity, ranging from fully parallel and communication-free (Independent) to fully coupled with peer consensus (Decentralized). These configurations allow us to test whether performance gains arise from \emph{agent coordination and specialization} or merely from increased compute through ensembling. \textcolor{black}{Our taxonomy covers coordination patterns common in LLM-based agentic systems, focusing specifically on \emph{communication topology}, one of several orthogonal MAS design dimensions including agent specialization \citep{hong2024metagpt}, memory architecture, and aggregation strategy. Classical coordination mechanisms such as blackboard systems assume structured message formats rather than natural language, limiting their direct applicability to LLM-based agents \citep{guo2024large, xi2025rise}.}

\textcolor{black}{We formally define the task-level error rate as $E = 1 - P$, where $P$ is the fraction of tasks successfully resolved. The \emph{task-level} error amplification factor $A_e^{\text{task}} = E_{\text{MAS}} / E_{\text{SAS}}$ quantifies the relative error rate of a multi-agent system compared to its single-agent baseline; $A_e^{\text{task}} > 1$ indicates that coordination introduces net errors, while $A_e^{\text{task}} < 1$ indicates net error suppression. We additionally define a \emph{trace-level} error amplification factor $A_e^{\text{trace}}$ that measures how much extra computational work arises from inter-agent coordination failures, estimated from execution-trace token analysis (see Section~\ref{subsec:coord_eff}). Both metrics consistently show that architectures with verification mechanisms contain errors more effectively than independent coordination, though they differ in absolute magnitude ($A_e^{\text{task}} \approx 1.1$--$1.3$ vs.\ $A_e^{\text{trace}} \approx 4$--$17$) because they capture complementary aspects of error dynamics.}


\subsection{Agentic Tasks and Benchmarks}

Following and extending the framework of \citet{zhu2025establishing}, we operationalize a task $T$ as \textbf{agentic} when optimal performance \emph{substantially} benefits from adaptive interaction. Formally, if $\tau = \{(a_t, o_t)\}_{t=0}^T$ represents an interaction trajectory, then:
\[
\frac{\max_{\pi} \mathbb{E}[R(\tau)] - \max_{g} \mathbb{E}[R(g(x))]}{\max_{\pi} \mathbb{E}[R(\tau)]} > \delta,
\]

where $\pi$ represents an interactive policy, $g$ represents any single-forward-pass function, $R$ measures task success, $\delta$ is a task-dependent threshold, and the expectation is over task instances $x$ and stochastic environment dynamics. This definition captures tasks where interaction provides meaningful advantage over the best possible single-shot approach.

The expected return of an optimal policy thus hinges on sequential observation–action feedback, requiring agents to gather information, plan, and revise hypotheses under partial observability. Building on the Agentic Benchmark Checklist \citep{zhu2025establishing}, we formalize three necessary properties for agentic benchmarks:

\begin{itemize}
    \item \textbf{Sequential Interdependence}: Later actions depend on earlier observations; a one-shot policy cannot achieve high reward.  
    \item \textbf{Partial Observability}: Critical state information is hidden and must be acquired through active querying or tool use.  
    \item \textbf{Adaptive Strategy Formation}: The policy must update internal beliefs based on new evidence obtained through interaction.
\end{itemize}

\textcolor{black}{\noindent Benchmarks lacking these conditions (e.g., GSM8K, MMLU) evaluate static reasoning rather than agentic capabilities. (We note that ``agentic'' is defined relative to current model capabilities: GSM8K could be posed as agentic by providing calculator tools, though current LLMs do not \emph{require} such scaffolding; conversely, tasks that are agentic today, such as SWE-Bench, may become solvable via single-shot inference as models improve. Our evaluation focuses on tasks that \emph{currently} require multi-step interaction for non-trivial performance.)}

\vspace{0.4em}
\paragraph{Why Environment Feedback Matters.}
Real-world deployments such as coding assistants, financial analysts, and embodied robots operate under uncertainty and non-stationarity.  
Tasks solvable by direct prompting measure linguistic knowledge, whereas agentic benchmarks evaluate the process of intelligence: exploration, adaptation, and coordination.  
Hence, our benchmarks are chosen such that (i) base LLMs perform poorly in single-shot mode, and (ii) non-trivial performance requires multi-step environment interaction.

\vspace{0.4em}
\paragraph{Benchmark Design Principles.}
Extending the framework proposed by \citet{zhu2025establishing}, we introduce additional criteria to isolate \emph{architectural effects}:

\begin{itemize}
    \item \textbf{Controlled Tool Interface:} identical tool APIs and observation structures for all architectures to eliminate confounds from external feedback quality.
    \item \textbf{Controlled for Parametric Knowledge:} within each model family, evaluation emphasizes adaptive reasoning over memorized facts. Cross-family comparisons (Section~\ref{sec:results}) account for inherent knowledge base differences through baseline normalization.
    \item \textbf{Action–Observation Loop Length:} each benchmark enforces non-trivial trajectory length $L>3$ to ensure sequential reasoning.
    \item \textbf{Comparative Normalization:} scores are normalized to the best single-agent baseline, measuring coordination gain or loss.
\end{itemize}


\begin{table}[t!]
\centering
\textcolor{black}{\caption{Six agentic benchmarks used for evaluation.}}
\adjustbox{width=\textwidth}{%
\begin{tabular}{lll}
\toprule
\textbf{Benchmark} & \textbf{Task} & \textbf{Evaluation Design} \\
\midrule
BrowseComp-Plus (2025) & Web Browsing / Information Retrieval & Multi-website Information Location \\
\rowcolor{lightgray}
Finance-Agent (2025) & Finance & Entry-level Analyst Task Performance \\
Plancraft (2024) & Agent Planning & Minecraft Environment Planning \\
\rowcolor{lightgray}
WorkBench (2024) & Planning / Tool Selection  & Common business activities  \\
\textcolor{black}{SWE-bench Verified (2024)} & \textcolor{black}{Software Engineering} & \textcolor{black}{GitHub Issue Resolution} \\
\rowcolor{lightgray}
\textcolor{black}{Terminal-Bench (2025)} & \textcolor{black}{CLI Task Execution} & \textcolor{black}{System Admin / Security / ML Tasks} \\
\bottomrule
\end{tabular}%
}
\end{table}